\section{Penanganan Kesalahan Sintaksis}
\label{sec:ch10-error}

\subsection{Jenis Kesalahan}

Kesalahan sintaksis terjadi ketika aliran token tidak sesuai dengan
tata bahasa. Contoh umum:
\begin{itemize}[leftmargin=*]
  \item titik koma hilang;
  \item kurung tidak seimbang;
  \item kata kunci yang salah;
  \item ekspresi tidak lengkap;
  \item deklarasi fungsi tanpa nama.
\end{itemize}

Kesalahan sintaksis berbeda dari kesalahan leksikal. Lexer hanya
melihat karakter, sedangkan parser melihat token dan struktur.
Karena itu, parser dapat memberikan pesan yang lebih kontekstual.

\subsection{Strategi Pemulihan}

Tiga strategi pemulihan yang umum:
\begin{enumerate}[label=\arabic*.,leftmargin=*]
  \item \textbf{panic mode}: lewati token sampai titik sinkronisasi;
  \item \textbf{phrase-level recovery}: perbaiki token lokal, misalnya
        sisipkan titik koma yang hilang;
  \item \textbf{error production}: tambahkan aturan yang mengenali
        kesalahan umum dan menghasilkan pesan khusus.
\end{enumerate}

\subsection{Mode Panik pada Predictive Parser}

Pada predictive parser, mode panik bekerja dengan himpunan token
penyelaras \(H\). Token pembatas seperti titik koma dan kurung tutup
dimasukkan ke \(H\), demikian pula token yang mengawali setiap
pernyataan. Jika \(M(A,a)\) bernilai error dan \(a \in H\), parser
membuang \(a\) dan token berikutnya sampai menemukan anggota \(H\),
lalu melanjutkan. Jika \(a \in \mathrm{FOLLOW}(A)\), entri tabel
dapat diubah menjadi \code{synch}: parser melepas \(A\) dari stack
dan melanjutkan tanpa membaca token.

Pemilihan \(H\) menentukan kualitas pemulihan. Jika \(H\) hanya berisi
\code{FOLLOW}, parser dapat melewati token yang mengawali pernyataan
berikutnya. Karena itu, token awal pernyataan harus dimasukkan ke
\(H\) secara eksplisit. Mode panik ini sederhana, tetapi kurang
efektif bila satu baris memuat beberapa kesalahan sekaligus.

\begin{figure}[H]
\centering
\begin{tikzpicture}[node distance=6mm and 9mm]
  \node[decision, text width=2.6cm] (err) {Kesalahan?};
  \node[phase, text width=2.8cm, below=of err] (report)
        {Laporkan lokasi};
  \node[phase, text width=2.8cm, below=of report] (sync)
        {Sinkronisasi};
  \node[phase, text width=2.8cm, right=of sync] (continue)
        {Lanjutkan parsing};
  \draw[arrow] (err) -- node[right, font=\small] {ya} (report);
  \draw[arrow] (report) -- (sync);
  \draw[arrow] (sync) -- (continue);
  \draw[arrow] (continue) |- (err);
\end{tikzpicture}
\caption{Alur panic mode pada parser}
\label{fig:ch10-panic}
\end{figure}

\subsection{Pesan Kesalahan yang Baik}

Pesan kesalahan sebaiknya memuat:
\begin{itemize}[leftmargin=*]
  \item lokasi: baris dan kolom;
  \item harapan: token apa yang diharapkan;
  \item temuan: token apa yang ditemukan;
  \item saran: perbaikan yang mungkin.
\end{itemize}

\begin{table}[H]
\centering
\caption{Contoh pesan kesalahan sintaksis}
\label{tab:ch10-error}
\begin{tabularx}{\textwidth}{@{}l X@{}}
\toprule
\textbf{Masukan} & \textbf{Pesan} \\
\midrule
\code{int a = 1} & \code{baris 1, kolom 10: diharapkan ';'
setelah deklarasi, ditemukan 'return'} \\
\code{if (a} & \code{baris 1, kolom 6: diharapkan ')' setelah kondisi} \\
\code{return ;} & \code{baris 1, kolom 8: ekspresi diharapkan setelah return} \\
\code{int (} & \code{baris 1, kolom 5: nama fungsi diharapkan} \\
\bottomrule
\end{tabularx}
\end{table}

\subsection{Batas Pemulihan}

Parser tidak dapat memperbaiki semua kesalahan. Jika masukan terlalu
rusak, parser sebaiknya berhenti setelah sejumlah kesalahan dan
melaporkan ringkasan. Strategi ini mencegah ledakan pesan dan
membantu pengguna fokus pada masalah utama.

\begin{infobox}
Kualitas pesan kesalahan sering menjadi pembeda antara compiler yang
menyenangkan dan yang menyulitkan. Investasi pada pesan yang baik
memberi manfaat besar bagi pengguna.
\end{infobox}
